\section{Background}
\label{sec:background}

\subsection{From static policies to context-adaptive access control}

Classical access control separates subjects, objects and operations. Role-based
access control (RBAC) assigns permissions to roles \cite{sandhu1996role};
attribute-based access control (ABAC) generalises this to attributes of the
subject, resource and environment \cite{hu2014guide}; usage control (UCON)
extends the decision to continuity, so that a permission can be revoked while an
action is in progress \cite{park2004ucon}. These models answer \emph{who may do
what}, but the security mechanisms they select are usually fixed by
configuration.

Risk-based access control makes the mechanism itself a function of a risk score
\cite{atlam2017overview,rath2017adaptive}. In enterprise settings, situational
awareness and socio-technical risk scores have been used to choose between
authentication factors \cite{lee2017situational,moura2017sotraace}. In IoT, CAAC
ties the decision to device context: battery, load, link quality, location and
neighbourhood behaviour \cite{genc2019context,arfaoui2019context}. Recent designs
add blockchain-backed continuous authentication \cite{hussainalnaji2022cab},
smart-contract policy evaluation \cite{merlec2024sc}, role--capability
combinations \cite{krishnasrija2025dynamic}, and real-time zero-trust postures
\cite{stodt2025beyond,federici2023zero}. Surveys document the breadth of the
design space and the recurring absence of an explicit adversary model
\cite{ravidas2019access,qiu2020survey,elhajj2019survey}. Context-aware computing
supplies the underlying abstraction: context is heterogeneous, noisy and
partially observable \cite{piccialli2021context}.

\subsection{Autonomic adaptation and its failure modes}

The MAPE-K loop---monitor, analyse, plan, execute over a shared knowledge
base---is the standard architecture for self-adaptive systems
\cite{kephart2003vision,cheng2009software}. Applied to security, it produces a
controller that observes context, updates a risk estimate, selects a policy and
enforces it \cite{elgendi2019protecting}. Two failure modes recur. First,
\emph{adaptation-induced downgrade}: because the loop maximises an expected
objective, a spoofed context that reduces the estimated risk also reduces the
selected security level. Second, \emph{thrashing}: when the context oscillates
near a decision boundary, the controller switches posture repeatedly, paying
reconfiguration and energy costs and creating a timing side channel
\cite{kephart2003vision}. Hysteresis and dwell timers are the usual mitigation,
but they interact with safety: a dwell timer that is too long can keep a
degraded posture after the risk has risen.

Runtime assurance takes a different route. The Simplex architecture keeps a
verified baseline controller and switches to it when the advanced controller is
untrusted \cite{sha2001simplicity}; shield synthesis and safety guards enforce a
temporal property at the actuator boundary \cite{wu2017safety,wu2019shield}; the
black-box Simplex architecture generalises the idea to learned components
\cite{mehmood2022black}; and recent work generates verified monitors from timed
automata \cite{cho2025timed}. The common principle is that the safety property
is enforced by a component that does not share the advanced controller's
uncertainty. Our design imports that principle into access control: I1/I2 are the
verified baseline, and the belief-driven scorer is the advanced controller that
may only choose within the baseline's envelope.

\subsection{Uncertainty, poisoning and Byzantine context}

Partially observable Markov decision processes provide the natural language for
acting under uncertain state \cite{kaelbling1998planning}, and hidden Markov
models provide a tractable forward filter when the state space is small
\cite{rabiner1989tutorial}. A belief is not automatically trustworthy: its
emission model can be wrong, and an adversary who understands the model can
choose observations that drive the posterior. False-data-injection and sensor
spoofing literature shows that stealthy attacks are designed exactly to remain
consistent with the defender's model \cite{barua2020special,li2019optimal}, while
summation and consistency checks detect attacks that break physical invariants
\cite{ye2020summation}. Byzantine fault tolerance supplies the aggregation-side
bound: with fewer than $N/2$ faulty nodes, a median or quorum can still be
informative \cite{lamport1982byzantine}, but the bound says nothing about
whether the \emph{honest} aggregate is being interpreted correctly.

Denial-of-sleep attacks add the reverse incentive: rather than hiding an attack,
the adversary makes the device believe it is under attack so that it spends
energy on authentication and inspection \cite{brownfield2005wireless,raymond2009effects}.
A design that only escalates on suspicion is therefore as exploitable as one
that only relaxes. The controller needs both directions and a defined failure
semantics.

\subsection{Calibration and lightweight enforcement}

Accuracy is the wrong metric for a risk estimate that feeds a decision. Strictly
proper scoring rules, especially the Brier score, reward honest uncertainty
\cite{gneiting2007strictly}, and modern calibration methods show that
high-accuracy classifiers can be badly miscalibrated \cite{guo2017calibration}.
We therefore report Brier score and expected calibration error alongside
detection AUROC, and we distinguish an estimator that is \emph{right} from one
that is \emph{confident and wrong}.

At the enforcement layer, authenticated encryption with associated data provides
confidentiality and integrity in one pass. Ascon is standardised for constrained
IoT devices \cite{dobraunig2021ascon}, and implementation studies report
microcontroller costs for it \cite{cagua2025implementation}. The action
space in this paper treats the choice of AEAD construction, authentication
strength, inspection depth, rate limiting and freshness as a joint posture; the
cryptographic primitives are assumed correct and are not re-proved.
